\section{模块 2 · 参考论文与技术路线：站在谁的肩膀上，为什么这么选}

做汇报时导师一定会问：``这方法是你发明的，还是参考了谁？''本模块把答案一次说清。
先给结论：\textbf{总体路线参考 Song \& Jung 的``栅格化 + 图像形态学''无监督工作流，
拆楼粘连用 Beucher \& Meyer (1993) 的经典分水岭算法，噪声清除的关键约定取自
荷兰 3D BAG 生产线（Peters et al.\ 2022）}；在此骨架上我们针对本区数据做了三处自己的改造
（不做闭运算、以高度面做分水岭地形、非对称碎片吸收——模块 3--4 详解）。

\subsection{选方法前先明确判断标准}

模块 1 说了本步的两个子任务（按栋分开、剔除假建筑）。挑方法时我们定了三条标准：

\begin{enumerate}
\item \textbf{可解释、可审计}：每个决定（这两栋为什么分开？那个实例为什么被删？）
      都要能指着中间产物讲清楚——这是本项目的验收习惯（每一步都在 CloudCompare
      里人工抽查，CloudCompare 是一款免费的点云查看软件）。
\item \textbf{不依赖额外训练数据}：语义分割已经花掉了标注预算，单体化最好``拿来就能用''。
\item \textbf{适配本区场景}：独栋住宅为主、楼间距小（1.5--2.5 米）、树多遮挡重。
\end{enumerate}

第 3 条不是一句空话，三个场景特点在方案里各有对应的招数：
楼间距小 $\to$ 分水岭按高度谷拆粘连（模块 4）；树冠被误标成建筑 $\to$
Gate A 与追溯修正逐实例清除（模块 5--6）；树遮挡造成屋顶空洞 $\to$
逐块补洞机制（模块 3）。当然它们各有局限，讲到时逐一说明。

\subsection{一个先决选择：要不要用官方``建筑底图''（footprint）}

在讲具体算法之前，有一个更根本的路线分岔。\textbf{footprint（建筑底图/轮廓）}
指政府地籍或测绘部门维护的每栋建筑的平面轮廓多边形——相当于官方已经画好的``每栋楼的地界''。

\begin{paperbox}[：两条路线]
\textbf{footprint-guided（有底图路线）}：把现成轮廓当输入，点云只负责``往轮廓里填高度''。
单体化问题被底图直接解决——一个轮廓就是一栋楼。这是\textbf{大型生产项目的常见路线}：
荷兰 3D BAG 项目以官方地籍轮廓 + LiDAR 重建了全国约 1{,}000 万栋建筑
（Peters, R.\ et al.\ 2022, \textit{PE\&RS} 88(3)。LoD2 指``带屋顶形状的三维建筑模型''
这一细节等级，是本管线的最终目标产品）；格拉斯哥的 LoD2 模型
（覆盖约 100 km\textsuperscript{2} 核心区，全域另有更粗的 LoD1）
也走有底图路线——值得注意的是，它的底图并非官方地籍，而是\textbf{从点云栅格掩膜
提取再正则化得到的轮廓}，配合开源重建工具 City3D（Huang et al.\ 2022）的
有底图模式使用。思想源头可追到 Vosselman \& Dijkman (2001) 用地面图辅助重建的经典工作。

\textbf{footprint-free（无底图路线）}：只用点云自己解决单体化——就是本项目的选择。
\end{paperbox}

为什么放着这条``省事的常见路线''不用？三个理由：
（a）\textbf{研究目标}本身就是检验``纯点云能走多远''——很多地区没有高质量底图，
方法要有普适性；
（b）官方底图有时效问题（新建、翻建、拆除不同步）；
（c）最重要的一条：本县（Franklin County）其实\textbf{有}官方 footprint，
但我们故意不让它进管线，而是留作最后第 9 步的\textbf{精度评价基准}
——同一份数据既当输入又当考卷，考试就没有意义了。

\subsection{主参考：Song \& Jung 的``栅格化 + 形态学''工作流}

\begin{paperbox}[：主参考文献]
Song, H., Jung, J.\ An Unsupervised, Open-Source Workflow for 2D and 3D
Building Mapping From Airborne LiDAR Data. \textit{IEEE JSTARS} 17
（2024 年正式收卷，2023 年在线发表；IEEE JSTARS 是美国电气电子工程师学会的
应用地球观测与遥感专题期刊）。
\end{paperbox}

这篇论文的核心思想只有一句话：\textbf{把三维点云``拍扁''成二维格子图，
然后用图像处理里成熟的工具来切建筑图斑}。``栅格化''``形态学''这些词模块 3 会从零讲；
这里先说清\textbf{为什么选它当骨架}：

\begin{enumerate}
\item \textbf{无监督}：整个流程不需要任何训练样本，拿到分类点云直接跑——满足标准 2。
\item \textbf{每一步都看得见}：中间产物是一张张图像，哪一步出错一眼可见——满足标准 1。
      后面会看到，我们正是靠``把中间图像画出来给用户检查''抓住了 v1 的大错误。
\item \textbf{计算便宜}：863 万个三维点的问题被压缩成一张几千$\times$几千的二维图像，
      个人电脑即可完成全区处理，无需服务器或显卡。
\item \textbf{天然衔接}：需要说明，原论文是从\textbf{原始} LiDAR 做起的完整工作流
      （含地面提取、建筑与树木的区分等环节）；我们\textbf{只借用其中
      ``栅格化 + 图像形态学制图''这一段}，并把自己第 1 步得到的建筑点作为这一段的输入
      ——相当于用深度学习语义分割替换了它的前半程，两者接口恰好吻合。
\end{enumerate}

要诚实说明的差别：Song \& Jung 的目标是\textbf{制图}（在地图上画出建筑图斑），
我们的目标是\textbf{实例}（每栋楼一个编号、还要把编号写回每一个三维点）。
所以我们沿用它``栅格化 + 形态学 + 连通域''的骨架，但在两处走得更远：
粘连楼的拆分（模块 4 的分水岭）和标签回写到点（模块 3 末尾）。

\subsection{拆粘连的关键工具：分水岭算法}

\begin{paperbox}[：经典算法文献]
Beucher, S., Meyer, F.\ (1993). The morphological approach to segmentation:
the watershed transformation. 收录于 \textit{Mathematical Morphology in Image
Processing}。——图像分割领域最经典的算法之一，30 多年来是``分开粘连目标''的标准工具
（显微镜下粘连的细胞、粘在一起的硬币，都是它的经典应用场景）。
\end{paperbox}

一句话直觉（模块 4 用图详细讲）：把图像看成\textbf{地形}，往里``灌水''，
水从不同的洼地涨起来，两片水相遇的地方\textbf{筑坝}——坝就是分割边界。
我们用它的\textbf{marker 控制版本}（marker-controlled watershed）：
先指定``水从哪几个源头涨起''，拆出来的块数就由源头数量控制，
噪声洼地不再各自成块——经典分水岭``一个噪声就多一块''的过分割毛病被压住了。
当然，责任随之转移到``源头找得准不准''上：找多了照样过分割、找少了照样欠分割，
所以模块 4 里``源头''的产生自带面积门槛把关。

\subsection{没有选的路线，以及为什么（导师大概率会追问）}

\textbf{（a）纯``连通域''为什么不够？}连通域（挨着的算一块，模块 3 详讲）
是最朴素的单体化。Yang et al.\ (2022, \textit{Int.\ J.\ Remote Sensing})
指出了这类做法的困难区间：\textbf{当楼间距逼近点间距时，简单的邻近分组容易失败、
相邻建筑难以分开}（他们因此提出``先把点组成小对象、再用对象间关系图合并成建筑''
的实例分割方法）。本区的账要算细一点：侧院 1.5--2.5 米、格距 0.5 米，
直接相除是 3--5 个格子宽的空隙；但屋檐每侧向外挑出 0.3--0.6 米，
两侧合计再吃掉 1--2 格，实际留在图上的空隙常常只有 2--4 格。
空隙本身不会让连通域出错——真正把空格``填上''的是三种东西：
落在缝里的零星误标点、伸进缝里的树冠、以及（v1 教训里的）闭运算。
缝越窄，越容易被填死——所以连通域只能当第一刀，粘连要靠分水岭补刀。

\textbf{（b）为什么不用更新的深度学习 / 图聚类实例分割？}同问题较新的代表工作之一是
Li, J.\ et al.\ (2026)：City-scale building instance segmentation via
structure-aware method（\textit{Int.\ J.\ Applied Earth Observation and
Geoinformation} 146:105086）。该文包含语义分割网络、新数据集等多项内容，
与我们对应的是其中的 \textbf{SI-DVDC 实例聚类}部分，思路是``密度可达的图聚类''：
把点（或小块）当作\textbf{节点}，空间上接近的节点之间连线，构成一张``图''
（网络结构意义上的图，不是图片）；能沿着足够密集的连线互相走到的节点，归为同一栋楼。
不选它当主线的理由：其一，它的强项是大城市复杂粘连建筑群，本区以独栋住宅为主，
用不上这份复杂度；其二，图聚类的中间状态难以像图像那样逐步画出来给人检查，
不符合我们``每步可审计''的工作方式；其三，引入新方法有调参与复现成本。
它被记为\textbf{候选兜底方案}（我们并未实际运行它，因此不称``对照''）
——如果形态学路线在某类场景失效，这是首选替代项。

\subsection{路线总览}

\begin{center}
\small
\begin{tabular}{p{3.1cm}p{3.6cm}p{3.9cm}p{3.6cm}}
\toprule
子问题 & 选用方法 & 依据文献 & 我们的改造 \\
\midrule
点云 $\to$ 可处理的图 & 0.5 m 占据栅格 & Song \& Jung 2024 & 每格 $\geq$2 点才算占据；\textbf{不做闭运算}（v1 教训，模块 3） \\
初步分块 & 8-连通域 & Song \& Jung 2024；困难区间见 Yang 2022 & 逐块补内部空洞（不向外桥接） \\
拆开粘连的楼 & marker 分水岭 & Beucher \& Meyer 1993 & 用 $-z_{\max}$ 高度面当地形；腐蚀取核当 marker（模块 4） \\
防过分割 & 碎片吸收 & —— （工程设计） & \textbf{非对称}规则：只许小并大（模块 4） \\
假建筑清除 & 离地高度 + 形状规则 + 人工判读 & 地面高度取``周边 4 m 内地面点第 5 百分位''借自 3D BAG（Peters 2022）；``桥必须在实例级清除''的机理见模块 5 & 3D BAG 依赖官方轮廓定``周边''，我们改用实例自身外扩 4 m 环；规则只``捞候选''，删除一律人工判读 + 原因码存档（模块 5） \\
候选兜底 & 图聚类实例分割 & Li et al.\ 2026 (SI-DVDC 部分) & 未运行，保留为替代路线 \\
\bottomrule
\end{tabular}
\end{center}

\begin{plainbox}[：模块 2 小结（汇报口径）]
一句话向导师交代：``我们走 footprint-free 路线（官方底图留作考卷），
主体沿用 Song \& Jung 的无监督栅格化工作流，粘连用经典 marker 分水岭拆分，
清噪沿用 3D BAG 生产线的离地高度约定；在闭运算、分水岭地形与碎片吸收三处
做了针对本区数据的改造，改造依据是 v1 版本的实测失败教训。''
\end{plainbox}
